\documentclass[crop,tikz]{standalone}
\usetikzlibrary{backgrounds}
\colorlet{blue}{cyan}
\tikzset{
  inverted/.style = {
    color=white,
    background rectangle/.style={fill},
    show background rectangle
  }
}

% Plots vector field and its magnitude of a finite solenoid of length
% h using the expressions from Martín-Luna et.al., "On the Magnetic
% Field of a Finite Solenoid", IEEE TRANSACTIONS ON MAGNETICS,
% VOL. 59, NO. 4, APRIL 2023

\usepackage{pgfplots}
\pgfplotsset{compat=1.18}
\usepgfplotslibrary{colormaps}
 
\begin{document}
\begin{tikzpicture}[inverted,inverted]
  % Note that the expressions are written in terms of the
  % dimensionless ratio r/r_0. The x-axis shows the ratio x/r_0 and
  % the y-axis shows the ration y/r_0. I.e. the borders of the
  % solenoid are at -1 and +1.
  \pgfmathsetmacro{\height}{2}; % ration h/r_0
  \pgfmathsetmacro{\xmin}{-3};
  \pgfmathsetmacro{\xmax}{3};
  \pgfmathsetmacro{\ymin}{-3};
  \pgfmathsetmacro{\ymax}{3};
\begin{axis}[inverted,
  xmin = {\xmin}, xmax = {\xmax},
  ymin = {\ymin}, ymax = {\ymax},
  axis equal image,
  view = {0}{90},
  hide axis,
  colormap={whiteblack}{rgb=(1,1,1) rgb=(0,0,0)},
  ]
  \addplot3 gnuplot [
    raw gnuplot,
    surf,
    shader=interp,
    no marks,
  ] {
set samples 101;
set isosamples 101;
set xrange [\xmin:\xmax];
set yrange [\ymin:\ymax];
set urange [\xmin:\xmax];
set vrange [\ymin:\ymax];
B0 = 4*pi;
h = \height;
k(r,z,s) = sqrt(4.*r/((1. + r)**2 + (z + 0.5*s*h)**2));
kp(r,z) = k(r,z,1.);
km(r,z) = k(r,z,-1.);
sigma(r) = 4.*r/(1. + r)**2;
Ipm(r,z,k) = -4./sqrt(r)*(EllipticE(k)/k - (1. - 0.5*k**2)*EllipticK(k)/k);
Jpm(r,z,k,s) = k/sqrt(r)*(z + 0.5*s*h)*((1. - r)/(1. + r)*EllipticPi(sigma(r),k) + EllipticK(k));
Ip(r,z) = Ipm(r,z,kp(r,z));
Im(r,z) = Ipm(r,z,km(r,z));
Jp(r,z) = Jpm(r,z,kp(r,z),1);
Jm(r,z) = Jpm(r,z,km(r,z),-1);
Bz0(r,z) = 0.5*B0*((z + 0.5*h)/sqrt(1 + (z + 0.5*h)**2) - (z - 0.5*h)/sqrt(1 + (z - 0.5*h)**2));
Br(r,z) = r == 0 ? 0 : B0/(4*pi)*(Im(r,z) - Ip(r,z));
Brz(r,z) = r == 0 ? Bz0(r,z) : B0/(4*pi)*(Jp(r,z) - Jm(r,z));
r(x,y,z) = sqrt(x**2 + y**2);
Bx(x,y,z) = r(x,y,z) == 0 ? 0 : Br(r(x,y,z),z)*cos(atan2(y,x));
By(x,y,z) = r(x,y,z) == 0 ? 0 : Br(r(x,y,z),z)*sin(atan2(y,x));
Bz(x,y,z) = Brz(r(x,y,z),z);
Ba(x,y,z) = sqrt(Br(r(x,y,z),z)**2 + Bz(x,y,z)**2);
splot "++" using 1:2:(Ba($1,0,$2))
  };
  % solenoid
  \draw[blue,very thick] (-1,{-\height/2}) -- (-1,{\height/2});
  \draw[blue,very thick] (1,{-\height/2}) -- (1,{\height/2});
\end{axis}
\end{tikzpicture}
\end{document}
